% !TeX program = xelatex
% Overleaf: Compiler = XeLaTeX, Main document = main.tex.
% Numeric tables are generated from saved JSON by build_tables.py.
\documentclass[11pt,a4paper]{article}
\usepackage{kotex}
\usepackage[margin=24mm]{geometry}
\usepackage{amsmath,amssymb,booktabs,array,tabularx,longtable}
\usepackage{graphicx,xcolor,setspace,enumitem}
\usepackage[authoryear,round]{natbib}
\usepackage{tikz,pgfplots}
\usetikzlibrary{arrows.meta,positioning}
\pgfplotsset{compat=1.18}
\usepackage{caption}
\usepackage[unicode,hidelinks]{hyperref}
\usepackage[section]{placeins}
\setstretch{1.25}
\setlength{\parindent}{1em}
\setlength{\parskip}{0.35em}
\setlength{\emergencystretch}{2em}
\widowpenalty=10000
\clubpenalty=10000
\displaywidowpenalty=10000
\captionsetup{font=small,labelfont=bf}
\renewcommand{\refname}{참고문헌}
\renewcommand{\tablename}{표}
\renewcommand{\figurename}{그림}
\newcommand{\checkneeded}[1]{\textcolor{red!65!black}{[확인 필요: #1]}}
\newcommand{\sourceproject}{\citep{project2026}}
\newcommand{\smalltable}[1]{\begingroup\small\setlength{\tabcolsep}{5pt}\input{tables/#1}\endgroup}
\title{스펙트로그램 특징 융합과 MiniROCKET을 이용한\\전기차 모터-감속기 경량 고장진단\\[4pt]
\large 차수 정규화와 Raspberry Pi 4 추론 성능 검증\\[8pt]
\normalsize Lightweight Fault Diagnosis of Electric Vehicle Motor-Reducers\\Using Spectrogram Feature Fusion and MiniROCKET}
\author{\checkneeded{저자명 및 저자 순서}\\\small\checkneeded{소속 및 교신저자}}
\date{}

\begin{document}
\maketitle
\begin{abstract}
전기차 구동계의 차량 내 고장진단에는 분류 성능과 제한된 계산 자원에서의 처리 가능성을 함께 고려해야 한다.
본 연구는 IONIQ, KONA, NIRO의 전류 및 두 진동 센서에서 생성된 STFT 스펙트로그램을 수치 특징 시계열로 변환하고, 다변량 MiniROCKET과 로지스틱 회귀를 결합한 경량 모터-감속기 고장진단 파이프라인을 구성하였다.
초기 12채널 융합을 출발점으로 컬러맵의 근사 역변환, 주파수별 기저 제거, 절대 주파수축의 분포 특징 및 rpm을 이용한 차수축 특징을 도입하였다.
주행 시각을 그룹으로 사용하는 3분할 평가 기록에서 최종 87채널 모델의 차종별 Macro-F1은 0.9861, 0.9682, 0.9644였으며, 차종 간 단순 평균은 0.9729였다.
동일 청크 구성을 사용한 51채널 모델 대비 DEMAG에서 NORMAL로의 오분류는 704건에서 433건으로 38.5\% 감소했으며, 양방향 혼동의 감소율은 27.5\%였다.
Raspberry Pi 4의 IONIQ 모델 단건 측정에서 예열 후 추론 지연의 중앙값은 455.04ms, 99백분위수는 457.30ms로 나타나, 10초 입력 창에 대응하는 계산 구간의 처리 가능성을 확인하였다.
다만 전체 입력으로 표준화와 MiniROCKET을 사전 적합한 평가 구조, 미확인 차수축 기준 회전수, STFT와 특징 추출을 제외한 지연 측정의 한계가 있으므로, 결과는 저장된 실험의 탐색적 성능과 부분 파이프라인의 엣지 실행 가능성으로 해석하였다.
\end{abstract}
\noindent\textbf{주제어:} 전기차, 모터-감속기, 고장진단, 다센서 융합, MiniROCKET, 차수 정규화, 엣지 추론

\section{서론}
\subsection{문제 제기 및 연구 배경}
전기차의 보급 확대는 구동계 상태를 지속적으로 판별하는 기술의 활용 범위를 넓히고 있다.
국제에너지기구는 2024년 세계 전기 승용차 판매가 1,700만 대를 넘었으며, 이 집계가 배터리 전기차와 플러그인 하이브리드를 포함한다고 보고하였다 \citep{iea2025}.
차량의 운행을 담당하는 모터와 감속기에서 발생하는 결함은 상태 감지와 대응이 필요한 대상으로, 전류와 진동을 이용한 진단이 연구되어 왔다 \citep{park2024}.
차량 내에서 이러한 진단을 수행하려면 분류 성능뿐 아니라 입력 처리, 추론 지연 및 메모리 사용을 함께 검토해야 한다.
특히 고해상도 스펙트로그램을 사용하는 접근에서는 이미지 표현과 계산량 사이의 관계가 설계 선택에 영향을 준다 \citep{park2024}.
이에 본 연구는 공개된 전기차 고장진단 자료로 경량 시계열 분류 구조를 설계하고, 저사양 단일 보드 컴퓨터에서 해당 구조의 추론 구간을 검증한다.

고장 상태를 분류하는 실험과 고장의 발생 시점을 예측하는 예지보전은 서로 구분할 필요가 있다.
본 프로젝트의 자료는 정상 및 네 결함 유형을 구분하는 라벨을 제공하므로, 직접 평가할 수 있는 과제는 다중 클래스 고장진단이다 \citep{aihub2022}.
현재 기록에는 실제 고장 발생 전후의 장기 추적 자료나 잔여수명 정답이 포함되어 있지 않다.
따라서 본 연구는 예지보전 시스템에 활용될 수 있는 진단 요소를 다루지만, 조기 탐지 시간이나 사고 예방 효과를 측정했다고 주장하지 않는다.
감자 결함의 미검출은 결함을 정상으로 판정하는 오류라는 점에서 별도로 검토한다.
구체적인 안전 위험 감소율이나 리콜 발생률은 프로젝트 자료에서 확인되지 않아 정량 근거로 사용하지 않는다.

\subsection{선행연구와 연구의 범위}
\citet{ku2023}는 진동 센서로 측정한 전기차 구동부 데이터를 임베디드 시스템에서 무선으로 서버에 전송하고, 데이터베이스와 머신러닝을 연계하는 상태 판별 구조를 제시하였다.
이 연구는 센서 수집, 상태 표시 및 서버 측 재학습을 연결한다는 점에서 진단 서비스의 시스템 구성을 설명한다.
그러나 제공된 1쪽 초록에는 분류 성능 수치, 데이터 분할 절차 또는 차량 내 분류기의 지연 측정이 제시되지 않는다.
따라서 해당 연구의 장점은 서비스 구조로 정리하되, 엣지 실행 성능에 대한 비교값은 확인할 수 없다.
\citet{jung2022}는 모터 테스트베드의 전류를 수정된 recurrence plot으로 이미지화하고, CNN 특징과 마할라노비스 거리를 이용한 건전성 지표를 제시하였다.
제공된 본문에서는 정상과 권선 단락의 심각도별 구분을 확인할 수 있으나, 본 연구의 3센서 융합이나 5개 구동계 클래스와 동일한 실험은 제시되지 않는다.

\citet{park2024}는 1280$\times$2000 스펙트로그램을 입력하는 DeepLabV3+ 기반 구조로 모터-감속기의 고장 분류와 결함 주파수 대역의 세그멘테이션을 함께 수행하였다.
이 연구는 AI-Hub의 전류 및 진동 자료를 사용하며, 분류 정확도 99.63\%와 클래스별 세그멘테이션 지표를 보고하였다.
그러나 제공된 2쪽 논문에는 Raspberry Pi와 같은 특정 엣지 장치에서의 지연 및 전력 측정이 제시되지 않는다.
이러한 보고 범위의 차이는 본 연구에서 하드웨어 실측을 수행하는 동기가 되지만, 선행 모델이 엣지에서 실행 불가능함을 의미하지는 않는다.
또한 본 연구의 출력은 청크 단위 클래스이므로, 결함 대역을 직접 표시하는 세그멘테이션 기능까지 대체했다고 볼 수 없다.
본 연구의 범위는 다센서 특징 시계열의 표현 개선, 경량 분류 구조 및 지정된 하드웨어에서의 추론 측정을 함께 검토하는 데 있다.

\subsection{연구 목적과 연구 질문}
본 연구의 목적은 차량 내 적용을 검토할 수 있는 모터-감속기 고장진단 파이프라인을 설계하고, 저장된 실험 기록을 통해 정확성과 계산 비용의 관계를 분석하는 것이다.
첫 번째 연구 질문은 MiniROCKET과 선형 분류기를 결합한 구조가 정상, 편심 두 단계, 감자 및 감속기 결함의 다섯 클래스를 어느 수준으로 구분하는가이다.
두 번째 연구 질문은 초기 다센서 융합과 이후의 절대축·차수축 특징 결합이 DEMAG와 NORMAL 사이의 혼동을 어떻게 변화시키는가이다.
세 번째 연구 질문은 해당 모델의 예열 후 계산 구간이 Raspberry Pi 4에서 10초 입력 창에 대응하는 처리 예산을 만족하는가이다.
이를 위해 초기 12채널 결과는 개발 경과로 제시하고, 최종 분석은 87채널 모델 및 동일 청크 구성을 사용하는 비교 모델을 중심으로 수행한다.
가설에 대한 통계적 검정은 기록에 없으므로, 각 연구 질문에는 관찰된 수치와 검증 범위를 명시하는 방식으로 답한다.

\section{이론적 배경}
\subsection{고장 유형과 시간-주파수 표현}
전류와 진동 신호는 모터 및 감속기의 운전 상태를 반영하는 측정 변수로 사용된다 \citep{aihub2022,park2024}.
공개 데이터의 모터-감속기 라벨은 NORMAL, ECC10, ECC20, DEMAG 및 REDUC의 다섯 범주로 구성된다.
ECC10과 ECC20은 데이터에서 편심 주의 및 편심 결함에 대응하고, DEMAG와 REDUC는 각각 감자 결함과 감속기 기어 결함에 대응한다 \citep{aihub2022}.
숫자 10과 20을 특정 물리 단위의 결함 크기로 해석하는 데 필요한 시험 조건은 본 연구에서 별도로 확인하지 않았다.
STFT는 시간 구간별 주파수 구성을 표현하므로, 비정상 운전 구간의 주파수 분포를 비교하는 입력으로 활용할 수 있다.
다만 본 연구는 이미 생성된 PNG를 사용하므로, STFT 창 길이와 중첩률 등 원신호 변환 설정은 \checkneeded{원본 STFT 설정}으로 남긴다.

이산 신호 $x[n]$의 STFT는 창 함수 $w$와 이동 간격 $h$를 사용하여 식~\eqref{eq:stft}와 같이 표현할 수 있다.
이 표현은 시간과 주파수의 두 축을 유지하지만, PNG는 원래 스펙트럼을 색으로 표시하고 양자화한 결과이다.
따라서 PNG의 RGB 밝기를 바로 신호 에너지로 사용하는 것은 색표의 스칼라 순서와 밝기가 일치한다는 추가 가정을 필요로 한다.
프로젝트에서는 이 가정이 성립하지 않는 사례와 대역 경계 특징의 포화를 확인하였다 \sourceproject.
또한 회전수에 따라 결함 대역의 bin 위치가 달라지는 관찰은 rpm에 따른 축 정렬을 검토하는 근거가 되었다.
이때 주파수 bin을 Hz로 바꾸는 보정 정보가 없으므로, 관찰된 bin/rpm 비율을 특정 고장의 물리적 특성 주파수 공식으로 단정하지 않는다.
\begin{equation}
 X(m,k)=\sum_n x[n]w[n-mh]\exp\left(-\mathrm{i}\frac{2\pi kn}{N}\right).
 \label{eq:stft}
\end{equation}

\subsection{이미지 분류와 특징 시계열 분류}
스펙트로그램 이미지 기반 CNN은 주파수 대역의 공간적 패턴을 학습하고, 세그멘테이션 구조를 통해 결함 위치도 예측할 수 있다 \citep{park2024}.
반면 ROCKET 계열은 시계열을 합성곱 통계 특징으로 변환한 뒤 선형 분류기를 학습하는 접근이다 \citep{minirocket2021}.
MiniROCKET은 고정된 커널 구조와 팽창 간격을 이용하여 반복적인 심층 신경망 학습의 부담을 줄인다.
본 프로젝트는 초기의 이미지 기반 후보와 수작업 요약 특징을 검토한 뒤, 특징 시계열에 MiniROCKET을 적용하는 구조로 전환하였다 \sourceproject.
다만 이 입력은 원시 센서 시계열이 아니라 PNG에서 추출한 대역 위치와 분포의 시계열이라는 점에서 원신호 직접 분류와 구분된다.
이미지의 전 영역을 유지하는 방식보다 표현 크기를 줄일 수 있지만, 선택되지 않은 공간 정보와 원신호의 위상 정보가 보존된다고 보장할 수는 없다.

MiniROCKET의 핵심 변환 통계는 합성곱 출력이 적합된 편향값을 초과하는 시간 비율로 설명할 수 있다 \citep{minirocket2021}.
본 연구에서는 다변량 구현으로 세 센서의 특징 채널을 함께 처리한다.
변환 후에는 로지스틱 회귀를 사용하여 다섯 클래스의 확률을 계산한다.
프롬프트에 기재된 RidgeClassifierCV는 가능한 비교 후보이지만, 최종 저장 모델과 학습 코드에서 확인되는 분류기는 로지스틱 회귀이다 \sourceproject.
따라서 본 논문은 수행되지 않은 Ridge 비교 성능을 제시하지 않는다.
이미지 모델 및 76차원 수작업 특징 기반 SVM·RF와의 동일 조건 비교 실험 역시 최종 평가 파일에서 확인되지 않아 방법 선택의 맥락으로만 다룬다.
\begin{equation}
 \phi_j(\mathbf{x})=\frac{1}{T_j}\sum_{t=1}^{T_j}
 \mathbb{I}\big[(\mathbf{x}*\mathbf{w}_j)_t>b_j\big].
 \label{eq:ppv}
\end{equation}

\subsection{분류 성능과 엣지 적합성의 측정}
Accuracy는 전체 청크 중 정답 비율을 나타내므로 시스템의 평균적인 분류 결과를 간단히 요약한다.
그러나 클래스의 비중이 다를 때에는 다수 클래스의 성능이 결과를 지배할 수 있다.
Macro-F1은 클래스별 F1을 동일 가중치로 평균하므로, 정상과 각 결함 유형의 균형 있는 성능을 확인하는 데 유용하다.
DEMAG 재현율은 실제 감자 결함 중 올바르게 검출된 비율을 나타내므로, 해당 클래스의 미검출을 별도로 평가할 수 있다.
본 연구는 폴드별 Macro-F1과 Accuracy의 평균·표준편차, 누적 혼동행렬에서 계산한 DEMAG 재현율을 구분하여 보고한다.
또한 엣지 적합성은 같은 입력 크기에서 예열 후 지연의 중앙값과 꼬리 백분위수를 기록하는 방식으로 평가한다 \sourceproject.
\begin{align}
 \mathrm{Accuracy}&=\frac{\sum_{c=1}^{C}TP_c}{n},\\
 F_{1,c}&=\frac{2TP_c}{2TP_c+FP_c+FN_c},\qquad
 \mathrm{Macro\!\text{-}\!F1}=\frac{1}{C}\sum_{c=1}^{C}F_{1,c},\\
 R_{\mathrm{DEMAG}}&=\frac{TP_{\mathrm{DEMAG}}}{TP_{\mathrm{DEMAG}}+FN_{\mathrm{DEMAG}}}.
\end{align}

\section{연구 방법}
\subsection{데이터 출처와 분석 대상}
연구 자료는 AI-Hub의 2022년 구축 데이터인 ``자율주행 고장진단 데이터''의 모터-감속기 부분이다 \citep{aihub2022}.
공식 소개에는 실제 환경에서 자체 수집한 자료로 설명되어 있으며, 모터-감속기 PNG와 대응 JSON이 각각 550,800건으로 기재되어 있다.
이는 제공 데이터 전체의 구축량이며 본 연구의 학습 청크 수를 의미하지 않는다.
프로젝트의 실험 정리에는 38,160개 Parquet와 센서 세 종류가 대응하는 12,720개 융합 구간이 기록되어 있다 \sourceproject.
최종 87채널 평가에는 클래스별 상한과 주행별 균등 추출을 적용한 46,660개 청크 및 1,176개 timestamp 그룹이 사용되었다.
연구자 확인에 따르면 최종 87채널 모델에는 AI-Hub 본데이터를 변환한 자료만 사용되었으며, 로컬 샘플 폴더의 합성·가공 자료는 포함되지 않았다.

원자료의 구축연도와 프로젝트 수행 시점은 구분하였다.
원자료는 2022년 구축 자료이며, 확보된 특징 개선 및 엣지 측정 기록은 2026년 7월 말부터 8월의 실험을 포함한다 \sourceproject.
연구자가 직접 센서를 설치하여 새로운 차량 데이터를 수집한 실험으로 서술하지 않는다.
원자료 수집의 정확한 시작일과 종료일은 제공 설명서에서 확인되지 않았으므로 \checkneeded{원자료 수집 시작·종료일}로 표시한다.
Current\_U는 모터 U상 전류, Vib\_Motor와 Vib\_TM은 각각 모터와 감속기의 수직 진동에 대응한다 \citep{aihub2022}.
차종별 모델은 따로 학습되므로 세 차종을 대상으로 하는 한 개의 공통 분류기를 학습한 것으로 해석하지 않는다.

\begin{table}[htbp]
\centering
\caption{최종 87채널 평가에 사용된 청크 및 timestamp 그룹 수. 클래스별 열은 청크 수가 아닌 주행 그룹 수이다.}
\label{tab:samples}
\smalltable{samples.tex}
\end{table}

\subsection{전처리와 특징 시계열 구성}
전처리 절차는 스펙트로그램 정렬, 근사 스칼라 복원, 특징 추출, 청크 및 센서 융합, 학습·엣지 측정의 다섯 단계로 수행되었다.
첫째, 차종·클래스·센서·측정 시각 및 분할 식별자가 대응하는 PNG 조각이 시간 순서로 연결되었다.
둘째, 색표 후보를 점검한 뒤 jet을 가정하는 LUT로 RGB 픽셀의 정규화 스칼라가 근사 복원되었다.
셋째, 주파수별 기저 제거를 적용한 위치 특징과 기저 제거 전 스칼라의 크기·분포 특징이 계산되었다.
넷째, 대응되는 세 센서가 고정 순서로 결합되고 길이 2,000의 비중첩 청크가 생성되었다.
다섯째, 차종별 MiniROCKET 및 로지스틱 회귀 모델이 적합되었으며, 저장 모델의 계산 구간이 Raspberry Pi 4에서 측정되었다 \sourceproject.

컬러맵 복원은 RGB 밝기와 색표 스칼라를 구분하기 위해 도입하였다.
코드는 jet 색표의 256개 기준색과 RGB 공간의 $64^3$ 복셀 중심 사이의 거리를 계산하고, 가장 가까운 기준색의 인덱스를 0과 1 사이로 정규화한다.
검사한 9개 PNG 중 7개에서는 jet이 후보 중 최적이었으나, 일부 파일에는 평탄 영역에서도 큰 색 잔차가 기록되어 있다 \sourceproject.
따라서 이 단계의 결과는 원신호의 전력이나 절대 에너지를 정확히 복원한 값으로 정의하지 않는다.
이하에서는 LUT 출력 $e_{f,t}$를 스펙트럼 강도 대용값이라고 부른다.
해당 값에는 원래 이미지의 정규화, 양자화 및 보간 효과가 남을 수 있으므로 물리 단위의 비교에는 추가 보정이 필요하다.

기본 네 특징은 peak\_freq\_bin, band\_start, band\_end 및 rms로 구성된다.
최종 추출기는 각 주파수 행의 시간 중앙값과 MAD를 이용해 기저를 제거하고, 그 결과의 전체 90백분위수를 임계값으로 사용한다.
피크 bin은 기저가 제거된 맵의 시간 열별 최대값 위치이며, 원 스펙트럼의 절대 최대 에너지 위치와 동일하다고 가정하지 않는다.
대역 시작과 끝은 임계값을 넘는 주파수 bin의 최저·최고 위치이므로, 하나의 연속된 결함 대역을 직접 추정하는 세그멘테이션 값도 아니다.
rms는 기저 제거 전 대용값의 제곱평균제곱근으로 계산되어 원시 전류나 가속도 신호의 RMS와 구분된다.
이러한 정의는 포화된 대역 경계와 밝기 기반 피크 위치 문제를 수정한 현행 코드에 따른 것이다 \sourceproject.
\begin{align}
 z_{f,t}&=\frac{e_{f,t}-\operatorname{median}_{t}(e_{f,t})}
 {1.4826\,\big[\operatorname{median}_{t}|e_{f,t}-\operatorname{median}_{t}(e_{f,t})|+10^{-6}\big]},\\
 \mathrm{rms}_t&=\sqrt{\frac{1}{H}\sum_{f=0}^{H-1}e_{f,t}^{2}}.
\end{align}

\subsection{절대축 및 차수축 확장 특징}
확장 특징은 주파수축을 균등 분할한 8개 대역의 비율과 centroid, entropy, rolloff 및 log\_flatness의 12개 변수로 구성된다.
대역 비율은 전체 스펙트럼 강도 대용값 중 해당 구간이 차지하는 비중을 나타내어, 위치 특징만으로 나타내기 어려운 분포 변화를 표현한다.
centroid와 entropy는 각각 분포의 중심과 퍼짐을 요약하며, rolloff는 누적 비율이 0.85에 도달하는 bin이다.
log\_flatness는 기하평균과 산술평균의 비를 로그 형태로 표현한다.
원시 에너지로 보정되지 않은 대용값을 사용하므로 이들 변수는 본 입력 표현 안에서의 분포 특징으로 해석한다.
프로젝트는 같은 12개 계산을 절대 주파수축과 rpm에 따른 재표본화 축에 적용하여 두 표현을 함께 사용하였다 \sourceproject.
\begin{align}
 p_{f,t}&=\frac{e_{f,t}}{\sum_{r=0}^{H-1}e_{r,t}+\epsilon},\qquad
 B_{j,t}=\sum_{f\in\mathcal{B}_j}p_{f,t},\\
 \mathrm{centroid}_t&=\sum_f f p_{f,t},\qquad
 \mathrm{entropy}_t=-\sum_f p_{f,t}\log(p_{f,t}+\epsilon),\\
 \mathrm{rolloff}_t&=\min\left\{r:\sum_{f=0}^{r}p_{f,t}\ge0.85\right\},\\
 \mathrm{log\_flatness}_t&=\frac{1}{H}\sum_f\log(e_{f,t}+\epsilon)
 -\log\left(\frac{1}{H}\sum_f e_{f,t}+\epsilon\right).
\end{align}

회전수 $r$의 조각에서는 낮은 주파수에서 시작하는 원본 bin $u$를 식~\eqref{eq:order}에 따라 선형 보간하여 차수축 대용 표현을 구성한다.
식에서 $o$는 0과 1 사이의 출력 축 좌표이며, $r_{\mathrm{ref}}$는 기준 회전수이다.
Hz 변환 정보가 없으므로 이 축의 값을 그대로 물리적인 무차원 차수 $f/(r/60)$와 동일시하지 않는다.
또한 $r<r_{\mathrm{ref}}$일 때에는 원본의 낮은 bin 범위를 출력 전체에 펼치므로, 전체 원주파수축을 손실 없이 유지하는 변환은 아니다.
실험 문서에는 본데이터의 회전수 범위 177--9,023rpm과 기준 9,100rpm이 기록되어 있지만, 현재 코드의 기본값은 5,700rpm으로 서로 다르다 \sourceproject.
최종 데이터 생성에 사용한 실제 기준값은 사용자가 현재 확인하기 어렵다고 답했으므로 \checkneeded{최종 87채널 특징 추출의 기준 rpm}으로 남기며, 그 설정의 재현성을 확보해야 한다.
\begin{equation}
 u(o;r)=o(H-1)\frac{r}{r_{\mathrm{ref}}}.
 \label{eq:order}
\end{equation}

최종 센서별 입력은 기본 4개, rpm 1개, 절대축 확장 12개 및 차수축 확장 12개를 합한 29개 특징이다.
이를 Current\_U, Vib\_Motor, Vib\_TM 순서로 결합하여 청크 하나를 $87\times2000$ 배열로 구성한다.
rpm은 세 센서에 반복하여 들어가므로 87개 채널이 모두 독립적인 측정 변수라는 뜻은 아니다.
세 센서의 길이가 다르면 코드는 가장 짧은 길이에 맞추고, 2,000행으로 나누어지지 않는 마지막 잔여 구간은 제외한다.
위치 특징 등 일부 열은 float32로 유지하고, 일부 비율·스칼라 열은 저장 시 float16을 사용한 뒤 학습에서 float32로 읽는다.
따라서 입력 표현의 경량성은 유지되지만, 최고 주파수 대역 비율의 언더플로우 가능성과 짧은 구간의 제외 효과도 함께 고려해야 한다 \sourceproject.

\begin{figure}[htbp]
\centering
\begin{tikzpicture}[node distance=5mm, box/.style={draw,rounded corners,align=center,text width=10.3cm,minimum height=8mm,font=\small},arr/.style={-{Latex[length=2mm]}}]
\node[box] (a) {전류 · 모터 진동 · 감속기 진동의 STFT 스펙트로그램 PNG};
\node[box,below=of a] (b) {jet LUT 근사 역변환 및 주파수별 기저 제거};
\node[box,below=of b] (c) {기본 4 + rpm 1 + 절대축 12 + 차수축 12 = 센서별 29특징};
\node[box,below=of c] (d) {3센서 early fusion: $87\times2000$ 청크};
\node[box,below=of d] (e) {채널 표준화 $\to$ MiniROCKET $\to$ 특징 표준화 $\to$ 로지스틱 회귀};
\node[box,below=of e] (f) {차종별 5클래스 판정 / Raspberry Pi 4 계산 구간 측정};
\draw[arr] (a)--(b);\draw[arr] (b)--(c);\draw[arr] (c)--(d);\draw[arr] (d)--(e);\draw[arr] (e)--(f);
\end{tikzpicture}
\caption{파이프라인 구성. 엣지 벤치마크는 이미 생성된 청크의 표준화부터 분류까지를 측정한다.}
\label{fig:pipeline}
\end{figure}

\subsection{모델 적합과 평가 절차}
차종별로 모델 한 개씩을 학습하고, 클래스당 최대 4,000개 청크를 주행 그룹별로 가능한 한 균등하게 선택하였다.
MiniROCKET의 요청 변환 특징 수는 10,000이며, 구현의 84개 기본 커널 조합 단위에 따라 실제 출력은 9,996개로 기록된다 \citep{sktime2026}.
본 논문에서 커널 수 스윕이라는 표현은 프로젝트의 명칭을 따르되, 표에는 실제 변환 특징 수를 표시한다.
데이터 및 변환의 난수 시드는 42이며, 최종 분류기는 반복 상한 300의 로지스틱 회귀를 사용하고 클래스 가중치는 지정하지 않는다.
분류기 앞의 StandardScaler는 각 폴드의 학습 특징에 적합되며, 최종 배포 모델은 전체 선택 청크로 다시 학습된다.
이 설정은 현행 학습 코드와 저장 리포트에 근거하며, 전체 파이프라인의 폴드 내부 적합 여부는 다음의 별도 제한과 함께 해석한다 \sourceproject.

교차검증에는 StratifiedGroupKFold를 사용하고, 최종 모델은 측정 시각 timestamp를 주행 그룹으로 사용한다.
이 구성은 동일 주행의 zsplit 또는 청크가 분류기의 학습과 검증 양쪽에 배치되는 것을 차단한다.
최종 자료의 클래스별 최소 timestamp 그룹 수는 48개로 기록되어, 초기 파일럿의 일부 클래스에서 그룹이 한 개뿐이던 조건과 다르다.
다만 현재 코드는 교차검증 분할 이전에 전체 선택 입력으로 채널 평균·표준편차를 계산하고 MiniROCKET의 편향을 적합한다.
그 뒤 변환된 특징에서 분류기와 특징 표준화만 폴드별로 다시 적합하므로, 검증 입력의 분포 정보가 전처리 및 변환 적합에 들어간다.
따라서 본 결과는 그룹 분할을 사용하는 탐색적 평가로 보고하며, 완전히 독립된 검증을 위해서는 채널 표준화와 MiniROCKET까지 폴드 내부에서 재적합해야 한다 \citep{sklearn2026}.

비교 실험은 rpm 추가, 온도 추가, 절대축 확장, 차수축 결합 및 변환 특징 수 축소를 포함한다.
15, 18, 48, 51, 81 및 87채널 실험은 차종별로 같은 청크 수와 주행 수를 기록하므로, 입력 구성 차이를 비교하는 중심 자료로 사용하였다.
동일한 표본 식별자가 실제로 사용되었는지는 원 인덱스와 분할 목록으로 추가 확인할 필요가 있다.
12채널 베이스라인은 차종당 20,000개 청크를 사용하여 최종 실험의 청크 구성과 다르므로, 12채널에서 87채널로의 변화 전체를 단일 통제 실험의 효과량으로 제시하지 않는다.
group\_key 비교 실험 역시 timestamp 실험과 청크 수가 달라, 수치 차이를 순수한 누수의 크기로 해석하지 않는다.
반복 시드 평가, 독립 테스트셋 및 통계적 유의성 검정은 저장 기록에 없어 수행된 것으로 서술하지 않는다 \sourceproject.

\subsection{엣지 추론 측정}
엣지 측정은 2026년 8월 23일 Raspberry Pi 4 Model B Rev 1.1의 64비트 aarch64 환경에서 수행된 JSON 기록을 사용하였다.
기록된 소프트웨어는 Python 3.13.5, NumPy 2.4.6, scikit-learn 1.7.2, sktime 1.1.0 및 Numba 0.67.0이다.
스레드 요청 수와 배치 크기의 조합은 $(1,1)$, $(4,1)$, $(4,8)$ 및 $(1,32)$이며, 단건은 30회, 배치 8은 10회, 배치 32는 5회 측정되었다.
JIT 및 캐시 예열 시간은 본 측정에서 분리하여 기록하였고, 각 배치의 전체 처리 시간과 청크당 처리 시간을 함께 계산하였다.
저장 모델의 측정 슬롯은 모두 IONIQ이므로, 세 차종의 엣지 지연을 각각 실측한 결과로 해석하지 않는다.
벤치마크는 실데이터 입력 파일이 없으면 난수 청크를 사용하는 코드이며, JSON에는 입력 파일 사용 여부가 기록되지 않아 \checkneeded{벤치마크 입력 종류 및 실행 명령}으로 남긴다 \sourceproject.

측정 구간은 준비된 청크의 채널 표준화, MiniROCKET 변환, 특징 표준화 및 확률 계산이다.
센서 수집, STFT, PNG 처리, LUT 역변환, 특징 추출 및 대시보드 통신은 포함하지 않는다.
따라서 계산 구간 지연은 전체 진단 응답시간보다 짧은 부분 측정으로 정의한다.
한 청크는 프로젝트에서 10초 PNG의 2,000개 시간 열에 대응한다고 설명되며, 이 주기의 계산 예산은 10,000ms로 설정한다.
초기 보고서의 PC 기준 목표 100ms와 이 입력 주기 기준은 서로 다르므로, Pi의 455ms가 초기 100ms 목표를 만족한다고 서술하지 않는다.
실시간 판단은 입력 창에 대응하는 처리량의 가능성을 검토하는 것이며, 안전 제어 시스템의 최악 실행시간 보증을 의미하지 않는다 \sourceproject.

\section{결과 및 논의}
\subsection{초기 센서 융합 결과}
초기 결과보고서는 개별 센서 특징의 평균 Macro-F1과 12채널 융합 결과를 비교하였다 \sourceproject.
보고서의 차종별 개별 센서 평균은 IONIQ 0.657, KONA 0.673, NIRO 0.628이며, 초기 융합 결과는 각각 0.871, 0.874, 0.826이다.
이 비교는 입력 단계에서 센서 정보를 함께 사용하는 방식의 유용성을 보여주는 개발 단계의 근거이다.
다만 보고서의 비교 표에는 개별 센서 점수의 평균이 기재되어 있으므로, 이를 확률 합산으로 구현한 late fusion의 실제 평가 점수와 동일하게 취급하지 않는다.
또한 초기와 후속 실험에서는 추출기, 그룹 정의 및 청크 구성이 바뀌었다.
따라서 이 수치들은 연구 질문 2의 초기 근거로 제시하되, 최종 특징 설계의 증가분을 산출하는 기준선으로 사용하지 않는다.

초기 혼동 분석에서 IONIQ의 Vib\_TM은 DEMAG를 NORMAL로 1,823건, NORMAL을 DEMAG로 1,782건 잘못 분류하였다.
같은 보고서의 12채널 융합에서는 두 방향의 오류가 각각 873건과 843건으로 줄었다 \sourceproject.
이 관찰은 여러 센서가 특정 경계에서 보완적인 정보를 제공할 수 있다는 해석과 부합한다.
그러나 평가 표본의 완전한 대응과 오류 수의 통계적 유의성은 추가 검증이 필요하다.
후속 분석에서는 RGB 밝기 기반 변환과 대역 경계 포화가 발견되어, 남은 혼동의 일부가 입력 표현에 의존할 가능성도 확인되었다.
따라서 초기 오류를 감자 고장의 물리적 분리 불가능성만으로 설명하는 것은 충분하지 않다.

\subsection{확장 특징과 차수축 결합의 효과}
표~\ref{tab:ablation}에 원본 평가 JSON의 차종별 폴드 평균을 정리하였다 \sourceproject.
rpm을 포함하는 15채널 구성의 차종 평균 Macro-F1은 0.8814였고, 절대축 확장 특징을 추가한 51채널 구성에서는 0.9635였다.
51채널에 차수축 특징을 더한 87채널에서는 평균이 0.9729로 증가하였다.
기본 특징과 rpm을 유지하면서 확장 특징을 차수축만으로 구성한 51채널 모델의 평균은 0.9562로, 두 축 결합보다 낮았다.
81채널 구성은 0.9723으로 87채널에 가까웠지만, 차종별 결과와 DEMAG 재현율에서는 완전히 같지 않았다.
이 결과는 평가에 사용한 입력 범위에서 분포 특징과 두 축의 결합이 유용했음을 보여주지만, 폴드 외부 변환 적합과 모델 선택에 사용한 평가의 반복을 고려해야 한다.

\begin{table}[htbp]
\centering
\caption{특징 구성별 Macro-F1. 각 차종 값은 3개 폴드 점수의 평균이며, 마지막 열은 차종 평균이다. 12채널은 선택 청크 수가 달라 직접적인 통제 비교에서 제외한다.}
\label{tab:ablation}
\smalltable{ablation.tex}
\end{table}

인버터 온도를 추가한 18채널 실험의 평균 Macro-F1은 0.8730으로, 15채널 구성보다 낮았다.
이를 모든 온도 정보가 진단에 무용하다는 결과로 일반화하지 않는다.
해당 실험에서 추가한 변수는 프로젝트 문서상 인버터 온도이며, 후속 CAN 자료의 모터 온도와는 구분된다 \sourceproject.
rpm 단독의 선형 진단 점수는 차종별로 0.208, 0.214, 0.238이었다.
낮은 단독 점수는 회전수 평균만으로 정답을 거의 결정하는 단순 지름길이 관찰되지 않았음을 시사한다.
그러나 비선형 rpm 규칙이나 다른 운전 조건과의 결합에 의한 지름길까지 배제하는 증거는 아니므로, 실제 운전 조건을 바꾼 검증이 필요하다.

\subsection{최종 분류 성능과 DEMAG 오류}
표~\ref{tab:final}에서 최종 모델의 IONIQ, KONA, NIRO Macro-F1은 각각 $0.9861\pm0.0008$, $0.9682\pm0.0053$, $0.9644\pm0.0074$ 수준으로 나타난다.
표의 정확한 표준편차는 저장 JSON에서 계산된 값을 반올림하여 제시하며, 이 값은 세 폴드의 산포이지 신뢰구간이 아니다.
차종 간 Macro-F1의 단순 평균은 0.9729이고, 누적 혼동행렬에서 계산한 DEMAG 재현율의 차종 평균은 0.9551이다.
Accuracy와 Macro-F1은 가까운 수준이지만, 특정 오류 경계의 남은 위험은 별도의 혼동행렬에서 확인해야 한다.
부록~\ref{app:cm}은 차종별 누적 혼동행렬을 제시한다.
이 결과는 연구 질문 1에 대한 탐색적 답변이며, 독립 테스트에서 같은 성능이 보장됨을 의미하지 않는다 \sourceproject.

\begin{table}[htbp]
\centering
\caption{최종 87채널 모델 성능. 평균$\pm$표준편차는 3개 폴드의 값이며, DEMAG 재현율은 누적 혼동행렬의 값이다.}
\label{tab:final}
\smalltable{final.tex}
\end{table}

51채널과 87채널을 비교하면 DEMAG 재현율은 세 차종 모두에서 증가하였다.
표~\ref{tab:demag}의 DEMAG에서 NORMAL로의 오류는 IONIQ 91건에서 42건, KONA 372건에서 199건, NIRO 241건에서 192건으로 변하였다.
이 방향의 오류 합계는 704건에서 433건으로 줄어 감소율이 38.5\%였다.
반대 방향까지 포함한 DEMAG와 NORMAL 사이의 양방향 오류는 1,462건에서 1,060건으로 감소하여 감소율은 27.5\%였다.
따라서 ``DEMAG와 NORMAL 혼동 38\% 감소''라는 표현은 전체 양방향 오류로 사용하지 않고 DEMAG 미검출 방향으로 한정한다.
이 차이는 분포 특징에 회전수에 따른 표현을 결합할 가능성을 보여주지만, 모든 차종에서 감자 오류가 해소되었다고 해석하지 않는다 \sourceproject.

\begin{table}[htbp]
\centering
\caption{51채널에서 87채널로의 DEMAG 재현율과 두 방향 오류 변화. $D$는 DEMAG, $N$은 NORMAL이며 pp는 퍼센트포인트이다.}
\label{tab:demag}
\resizebox{\linewidth}{!}{\smalltable{demag.tex}}
\end{table}

표현 개선의 결과는 DEMAG와 NORMAL의 분리가 입력 표현에 민감하다는 해석을 지지한다.
그러나 LUT 보정, 기저 제거, 확장 특징 및 차수축을 함께 검토한 개발 이력만으로 한 단계가 오류의 유일한 주원인이라고 확정할 수는 없다.
실제 프로젝트 기록에서도 LUT 보정만으로 모든 차종·클래스의 분리도가 일관되게 개선되지는 않았다 \sourceproject.
이는 잘못된 색 변환으로도 일부 클래스가 일관된 패턴을 보일 수 있으며, 물리적으로 더 타당한 표현과 분류 성능이 반드시 일대일로 연결되지는 않음을 시사한다.
절대축과 차수축의 결합은 한 표현이 잃는 정보를 다른 표현이 보완하는 후보 설계로 해석할 수 있다.
다만 채널 수 증가 자체와 표현의 상보성을 엄밀히 구분하려면, 중복 특징을 추가한 같은 채널 수의 대조군도 필요하다.

\subsection{변환 특징 수와 경량화의 절충}
87채널을 유지하면서 실제 변환 특징 수를 9,996, 4,956, 2,436 및 1,176으로 변경한 결과를 표~\ref{tab:kernels}에 제시한다.
Macro-F1의 차종 평균은 각각 0.9729, 0.9707, 0.9691 및 0.9587이었다.
기록된 차종별 모델 직렬화 크기는 기본 모델 약 0.69MB에서 가장 작은 구성 약 0.10MB까지 감소하였다 \sourceproject.
그러나 Pi 단건 지연은 9,996개에서 455.04ms, 2,436개에서 411.96ms로, 특징 수의 감소 비율만큼 줄지는 않았다.
1,176개 구성의 226.92ms는 더 큰 지연 감소를 보였지만, DEMAG 평균 재현율도 감소하였다.
따라서 이 측정 조건에서는 출력 특징 수 감소가 저장 크기에는 효과적이지만, 정확성과 추론 지연의 절충은 비례적이지 않았다.

\begin{table}[htbp]
\centering
\caption{87채널의 변환 특징 수 스윕. 분류 점수는 세 차종 평균, 모델 크기는 학습 리포트의 차종별 직렬화 크기, 지연은 Pi의 IONIQ·1스레드·배치 1 조건이다.}
\label{tab:kernels}
\resizebox{\linewidth}{!}{\smalltable{kernels.tex}}
\end{table}

\begin{figure}[htbp]
\centering
\begin{tikzpicture}
\begin{axis}[width=0.82\linewidth,height=6cm,xlabel={Pi 단건 지연 (ms)},ylabel={Macro-F1 (차종 평균)},xmin=200,xmax=490,ymin=0.95,ymax=0.98,grid=major,scaled y ticks=false,yticklabel style={/pgf/number format/fixed,/pgf/number format/precision=3}]
\addplot+[only marks,mark=*,mark size=2.5pt,point meta=explicit symbolic,nodes near coords,every node near coord/.append style={font=\scriptsize,anchor=south east}] table[x=latency,y=f1,meta=features]{tables/kernel_curve.dat};
\end{axis}
\end{tikzpicture}
\caption{변환 특징 수와 성능·지연의 관계. 점의 숫자는 실제 변환 특징 수이며, 분류 점수와 Pi 측정 슬롯의 범위가 다름에 유의한다.}
\label{fig:tradeoff}
\end{figure}

9,996개와 2,436개 구성의 지연 차이는 약 9.5\%로, 네 배가량의 출력 특징 수 차이보다 작다.
이 현상은 MiniROCKET의 기본 커널, 팽창 간격 및 입력 순회 비용을 함께 고려해야 함을 시사한다.
그러나 소스 수준의 연산 횟수 계측 없이 팽창 간격 상한만을 유일한 원인으로 확정하지 않는다.
또한 4,956개 구성은 2,436개보다 느리지만 Macro-F1이 높으므로, 모든 성능 축에서 완전히 지배당하는 대안이라고 표현하지 않는다.
입력 주기의 계산 예산을 기준으로 할 때 기본 모델도 충분히 짧은 측정 지연을 보여, 프로젝트는 분류 성능을 우선한 87채널·9,996개 구성을 선택하였다.
이 선택은 현재 비교 모델과 측정 조건 내의 판단이며, 최소 전력이나 최적 비용을 달성했다는 의미는 아니다 \sourceproject.

\subsection{Raspberry Pi 4의 추론 지연}
표~\ref{tab:latency}에서 기본 모델의 청크당 중앙 지연은 측정 조건에 따라 443.01--455.04ms였다.
1스레드·배치 1의 중앙값은 455.04ms이며, p95는 455.95ms, p99는 457.30ms, 관찰 최대값은 457.78ms였다.
10,000ms의 입력 창 예산을 중앙값으로 나눈 비율은 약 22.0으로, 준비된 청크에 대한 계산 구간에서는 입력 속도를 따라갈 가능성이 확인된다.
그러나 이 비율은 미측정 STFT와 특징 추출 비용까지 포함한 전체 시스템의 여유율이 아니다.
입력 창 10초를 모두 수집한 뒤 분류한다면 첫 결과까지의 시간에는 수집 창과 특징 생성 시간이 추가된다.
따라서 연구 질문 3에 대한 답은 예열 후 분류 계산 구간의 처리 가능성으로 한정하며, 종단 간 실시간성을 검증한 것으로 확대하지 않는다 \sourceproject.

\begin{table}[htbp]
\centering
\caption{Pi의 청크당 p50 지연(ms). $t$는 요청 스레드 수, $b$는 배치 크기이다. 배치 8과 32의 값은 전체 배치 지연을 배치 크기로 나눈 처리량 환산값이다.}
\label{tab:latency}
\smalltable{latency.tex}
\end{table}

기본 모델의 1스레드 단건 단계별 중앙 지연은 채널 표준화 1.351ms, MiniROCKET 452.08ms 및 분류 단계 1.568ms였다.
각 단계 중앙값의 합은 전체 중앙값과 반드시 같지 않지만, 변환 단계가 전체 중앙값의 약 99.35\%에 해당하는 규모임을 확인할 수 있다.
요청 스레드 수를 1에서 4로 바꾼 단건 지연은 455.04ms에서 454.18ms로 거의 변하지 않았다.
이 결과만으로 Pi에서 병렬화가 불가능하거나 정확히 CPU 코어 한 개만 점유한다고 단정할 수는 없다.
배치 32의 청크당 중앙 지연은 448.38ms였지만, 해당 배치 전체의 중앙 지연은 14,348.32ms이므로 단건 응답시간과 구분해야 한다.
따라서 온라인 운용에서는 배치 처리량의 개선보다 단건 측정과 입력 수집 지연을 함께 검토하는 것이 타당하다 \sourceproject.

기본 모델의 단건 예열 시간은 1.37초로 기록되었고, 배치 32에서는 42.96초였다.
이 예열 측정은 JIT 및 여러 예열 호출을 포함하므로, 모든 새 프로세스에서 동일한 냉시작 시간으로 보장되는 값은 아니다.
단건 기록에서 네 모델 레지스트리의 적재 시간은 0.25초, 적재 후 프로세스 RSS는 193.3MB였다.
해당 메모리는 기본 모델 하나의 사용량이나 프로세스의 최대 메모리로 해석하지 않는다.
네 측정 파일에 기록된 시작·종료 온도는 37.4--43.8\,${}^{\circ}$C 범위였고, 해당 시점의 스로틀 플래그는 모두 \texttt{0x0}, 클럭은 1,500MHz였다.
이 기록은 샘플링 시점의 상태를 확인해 주지만, 전체 실행 동안 무스로틀 상태였다는 연속 관측이나 전력 소비 측정까지 제공하지는 않는다 \sourceproject.

\subsection{선행연구와의 비교 및 연구 질문에 대한 답변}
표~\ref{tab:prior}는 선행연구와 본 연구의 과제 및 측정 범위를 비교한다.
박병준·한동석의 분류 정확도 99.63\%와 본 연구의 평균 Macro-F1 0.9729는 지표와 평가 절차가 달라 직접적인 우열 비교에 사용할 수 없다 \citep{park2024}.
동일한 공개 데이터 이름을 사용하더라도 학습 표본, 분할, 전처리 및 세그멘테이션 출력의 차이를 통제해야 공정한 비교가 가능하다.
본 연구의 차별성은 최고 점수의 주장보다 특징 표현을 수정한 과정과 특정 엣지 하드웨어에서의 추론 기록을 함께 제시한 데 있다.
또한 구교문의 서버 기반 구조 및 정원호의 건전성 지표는 각각 서비스 갱신과 심각도 표현의 장점을 제공하므로, 본 연구의 청크 분류와 대체 관계만으로 평가하지 않는다.
향후 동일 분할에서 이미지 및 특징 시계열 모델을 다시 비교해야 정확도·지연·전력의 실제 절충을 판단할 수 있다.

\begin{table}[htbp]
\centering\small
\caption{선행연구와의 비교. ``미제시''는 제공된 논문 범위에서 확인되지 않았음을 의미한다.}
\label{tab:prior}
\begin{tabularx}{\linewidth}{>{\raggedright\arraybackslash}p{25mm}*{3}{>{\raggedright\arraybackslash}X}}
\toprule
연구 & 입력 및 주요 과제 & 보고된 결과 & 엣지 측정 범위\\
\midrule
구교문 외(2023) & 진동 수집·서버 상태 판별 & 정량 분류 성능 미제시 & 장치별 추론 지연 미제시\\
정원호 외(2022) & 전류 RP·권선 단락 심각도 & CNN 특징 기반 건전성 지표 & 장치별 추론 지연 미제시\\
박병준·한동석(2024) & STFT 이미지·분류 및 세그멘테이션 & Accuracy 99.63\%; 클래스별 분할 지표 & Pi 지연·전력 미제시\\
본 연구 & 3센서 특징 시계열·5클래스 분류 & 평균 Macro-F1 0.9729 & Pi IONIQ 계산 구간 p50 455.04ms\\
\bottomrule
\end{tabularx}
\end{table}

연구 질문 1에 대해서는 최종 입력 표현에서 높은 탐색적 분류 점수가 관찰되었으나, 독립 검증 점수는 아직 확보되지 않았다고 답할 수 있다.
연구 질문 2에 대해서는 초기 센서 융합의 개선 경향과 후속 절대축·차수축 결합의 DEMAG 오류 감소가 관찰되었다.
다만 초기 비교와 후속 특징 비교는 서로 다른 평가 단계이며, 같은 효과량으로 합산하지 않는다.
연구 질문 3에 대해서는 Pi의 준비된 청크 계산 구간이 10초 입력 창의 시간 예산보다 짧았다고 답할 수 있다.
STFT부터 최종 출력까지의 전체 지연과 차량 내 안전 요구는 별도 검증 대상으로 남는다.
유의성 검정이나 사전 명시된 통계 가설이 없으므로, 세 연구 질문을 모두 통계적으로 채택된 가설로 표현하지 않는다.

\subsection{한계와 후속 연구}
첫 번째 한계는 그룹 분할과 별개로 전체 입력이 표준화와 MiniROCKET 적합에 사용된 평가 구조이다.
검증 입력의 분포가 변환 설정에 반영될 수 있으므로, 현재 점수의 낙관 편향 크기는 폴드 내부 재적합 없이는 산출할 수 없다 \citep{sklearn2026}.
다수의 특징 및 채널 구성을 같은 평가에 반복 비교한 점도 모델 선택 편향의 가능성을 남긴다.
이를 보완하려면 외부 그룹 테스트를 보류하고, 학습 내부에서만 특징 선택과 변환 적합을 수행하는 중첩 검증을 구성해야 한다.
둘째, 차종별 모델을 따로 적합하였으므로 미학습 차종의 일반화 성능은 평가되지 않았다.
후속 연구에서는 세 차종 중 한 차종을 제외하여 학습하는 leave-one-vehicle-out 실험과 날짜·장치 단위 외부 검증을 함께 수행할 필요가 있다.

셋째, 최종 입력은 AI-Hub 본데이터의 변환 자료로 확인되었으나, 실제 도로의 자연 발생 고장에 대한 일반화를 판단할 구체적인 결함 재현 조건은 추가 확인이 필요하다.
LUT 출력은 정규화 강도 대용값이며, 일부 색 잔차와 양자화 영향을 포함하므로 물리적 에너지의 완전한 복원으로 해석할 수 없다.
주파수 bin에서 Hz로의 보정과 원시 신호를 확보하면, 색표 근사 복원을 거치지 않는 입력과 직접 비교할 수 있다.
차수축의 실제 기준 rpm이 확인되지 않은 점 역시 동일 모델을 재생성하는 데 중요한 제한이다.
전체 구간의 중앙값과 백분위수를 이용하는 특징 추출은 미래 시간 정보를 사용할 수 있으므로, 실제 온라인 처리에는 과거 입력만 사용하는 인과적 추출 절차가 필요하다.
따라서 원시 신호와 최종 특징 사이의 손실, 실시간 처리의 인과성 및 설정 이력을 함께 검증해야 한다 \sourceproject.

넷째, 토크와 모터 온도 등의 CAN 참조 자료는 후속 전처리 문서에 존재하지만 최종 87채널 분류기의 추가 입력으로 검증되지 않았다.
특히 참조 토크와 실제 토크가 다수 파일에서 같은 값으로 기록된 문제가 있어, 토크 추종 오차를 유효한 특징으로 바로 사용할 수 없다 \sourceproject.
후속 연구는 원 자료의 신호 의미와 시간 정렬을 확인한 뒤, 토크·rpm·온도를 결합하는 조건부 진단과 운전 조건별 성능 분석을 수행해야 한다.
다섯째, Pi 측정은 한 차종 슬롯과 제한된 반복 수를 사용하였고, 실제 입력 사용 여부도 JSON만으로 확인되지 않는다.
STFT·특징 추출 포함 지연, 냉시작, 장시간 온도·CPU 사용 및 전력 계측을 추가해야 실제 배포 비용을 판단할 수 있다.
이 검증을 통해 경량 모델의 계산 가능성뿐 아니라 새로운 운전 조건에서의 진단 신뢰성을 평가할 수 있다.

\section{결론}
본 연구는 전기차 모터-감속기 스펙트로그램을 다센서 특징 시계열로 구성하고, 다변량 MiniROCKET과 로지스틱 회귀를 결합한 경량 진단 구조를 분석하였다.
입력 표현의 보정과 절대축·차수축 특징 결합을 적용한 87채널 모델에서 저장된 탐색적 평가의 차종 평균 Macro-F1은 0.9729였다.
동일 청크 수를 사용하는 51채널 대비 DEMAG에서 NORMAL로의 오류는 38.5\%, 두 클래스 사이의 양방향 오류는 27.5\% 감소하였다.
Pi의 IONIQ 단건 예열 후 계산 구간은 중앙값 455.04ms로 측정되어, 준비된 10초 입력 창에 대응하는 분류 처리의 가능성을 확인하였다.
이 성과는 표현 방식과 계산 구조를 함께 검토하는 접근의 유용성을 보여주지만, 폴드 외부 변환 적합과 미측정 전처리 때문에 독립 일반화와 종단 간 실시간성의 증거로 확대할 수 없다.
후속 연구는 폴드 내부 전체 파이프라인 적합, 차종·운전 조건 간 외부 검증 및 센서 입력부터 판정까지의 지연·전력 측정을 통해 실제 차량 적용 가능성을 검증해야 한다.

\begin{thebibliography}{99}
\raggedright
\bibitem[IEA(2025)]{iea2025}
International Energy Agency. (2025). Global EV Outlook 2025: Trends in electric car markets.
\url{https://www.iea.org/reports/global-ev-outlook-2025/trends-in-electric-car-markets-2}.

\bibitem[구교문 외(2023)]{ku2023}
구교문, 김효영, 김기현, 심재홍. (2023). 머신러닝을 이용한 전기 자동차 구동부의 상태 판별 시스템 개발.
한국생산제조학회 추계학술대회, 256.

\bibitem[정원호 외(2022)]{jung2022}
정원호, 임윤섭, 정성진, 박용화. (2022). 전류 데이터 이미지화 및 반지도학습을 이용한 전동기 고장 진단.
KSME\_DC22\_Thu\_18, 59--60. \checkneeded{학술대회 공식 서지명}.

\bibitem[박병준·한동석(2024)]{park2024}
박병준, 한동석. (2024). 이미지 세그멘테이션 모델을 활용한 실시간 전기차 모터-감속기 고장 진단법.
한국통신학회 하계종합학술발표회, 362--363.

\bibitem[Dempster et~al.(2021)]{minirocket2021}
Dempster, A., Schmidt, D. F., \& Webb, G. I. (2021).
MiniRocket: A very fast (almost) deterministic transform for time series classification.
Proceedings of the 27th ACM SIGKDD Conference on Knowledge Discovery \& Data Mining, 248--257.
\url{https://doi.org/10.1145/3447548.3467231}.

\bibitem[AI-Hub(2022)]{aihub2022}
AI-Hub. (2022). 자율주행 고장진단 데이터 [2022년 구축 데이터 및 데이터 설명서 v1.3].
\url{https://aihub.or.kr/aihubdata/data/view.do?aihubDataSe=data&currMenu=115&dataSetSn=71347&topMenu=100}.
온라인 소개 확인일: 2026년 10월 6일.

\bibitem[sktime(2026)]{sktime2026}
sktime. (2026). MiniRocketMultivariate API documentation [온라인 문서].
\url{https://www.sktime.net/docs/api-reference/sktimetransformationsrocketminirocketmultivariate/}.
확인일: 2026년 10월 6일. 실제 실험 버전은 본문에 별도 명시.

\bibitem[scikit-learn(2026)]{sklearn2026}
scikit-learn. (2026). Common pitfalls and recommended practices: Data leakage [온라인 문서].
\url{https://scikit-learn.org/stable/common_pitfalls.html}.
확인일: 2026년 10월 6일.

\bibitem[프로젝트 실험기록(2026)]{project2026}
프로젝트 실험기록. (2026). 모터-감속기 경량 고장진단 프로젝트 자료 [미간행 내부 자료].
프로젝트결과보고서; 발표자료; 특징 추출 및 학습 v3 코드; 구성별 평가 JSON;
엣지벤치 result\_t1, result\_t4, result\_t4\_b8, result\_t1\_b32 JSON.
세부 파일과 SHA-256은 동봉한 \texttt{source\_manifest.json}에 기재.
\end{thebibliography}

\clearpage
\appendix
\section{최종 모델의 누적 혼동행렬}\label{app:cm}
각 행렬은 세 검증 폴드의 혼동행렬을 합한 값이며, 행이 실제 클래스이고 열이 예측 클래스이다.
행 합은 해당 클래스의 청크 수이므로 주행 그룹 수와 구분한다.
평균 Macro-F1은 폴드별 점수의 평균이므로, 누적 행렬에서 다시 계산한 Macro-F1과 일치하지 않을 수 있다.
DEMAG 재현율과 방향별 오류 수는 이 누적 행렬에서 산출하였다.
원본 파일 경로와 수치 대응은 동봉한 \texttt{verified\_metrics.csv}에서 확인할 수 있다.

\begin{table}[htbp]\centering\caption{IONIQ: 최종 87채널 누적 혼동행렬.}\smalltable{cm_ioniq.tex}\end{table}
\begin{table}[htbp]\centering\caption{KONA: 최종 87채널 누적 혼동행렬.}\smalltable{cm_kona.tex}\end{table}
\begin{table}[htbp]\centering\caption{NIRO: 최종 87채널 누적 혼동행렬.}\smalltable{cm_niro.tex}\end{table}

\clearpage
\section{최종 확인이 필요한 재현 항목}
\begin{longtable}{p{35mm}p{110mm}}
\toprule 항목 & 확인 내용\\\midrule
저자 및 투고 형식 & 저자명·순서·소속·교신저자 및 목표 학회 양식. 현재는 일반 한국어 학술논문 형식이다.\\
데이터 계보 & 최종 학습은 AI-Hub 본데이터 변환 자료만 사용한 것으로 연구자 확인. 정확한 변환 실행 이력·원본 파일 인덱스 및 결함 재현 조건은 추가 확보 필요.\\
수집 및 STFT & 원자료의 수집 시작·종료일, 원신호 표본율, STFT 창·이동 간격, 색표 정규화와 Hz 축 보정.\\
차수축 기준값 & 실험 문서 9,100rpm과 현행 코드 5,700rpm의 차이. 사용자 답변: 현재 확인하기 어려움.\\
검증 재실행 & 폴드 내부 채널 표준화·MiniROCKET 적합 및 보류 테스트를 적용한 일반화 점수. 현재 점수로 대체하지 않는다.\\
벤치 입력·명령 & 실데이터/난수 사용 여부, 예열 횟수와 전체 실행 명령. JSON에는 입력 이력이 없다.\\
종단 간 성능 & 센서 수집부터 STFT·특징 추출·분류까지의 지연, 냉시작 및 실제 소비전력.\\
문헌 서지 & 정원호 외(2022)의 공식 학술대회명. 제공 PDF의 DC22 코드와 페이지를 우선 기재하였다.\\
\bottomrule
\end{longtable}
\end{document}
